% NeurIPS 2026 submission scaffold.
\documentclass{article}

\usepackage[preprint]{neurips_2026}
\usepackage[utf8]{inputenc}
\usepackage[T1]{fontenc}
\usepackage{hyperref}
\usepackage{url}
\usepackage{booktabs}
\usepackage{amsfonts}
\usepackage{amsmath}
\usepackage{amssymb}
\usepackage{nicefrac}
\usepackage{microtype}
\usepackage{xcolor}

\title{Context is the New Weight:\\
       When In-Context Information is Functionally Equivalent to Weight Updates}

\author{%
  Zizhao Hu \\
  Department of Computer Science \\
  University of Southern California \\
  \texttt{zizhaohu@usc.edu} \\
}

\begin{document}

\maketitle

\begin{abstract}
We ask when a context change applied to a fixed model produces \emph{exactly}
the same input-output behavior as a corresponding weight update on the same
model with empty context. We characterize a regime in which this equivalence
holds tightly, derive the context-to-weight map, and study its consequences
for fine-tuning, knowledge editing, and unlearning.
\end{abstract}

\section{Introduction}
\label{sec:intro}

% TODO

\section{Setup}
\label{sec:setup}

% TODO: define equivalence, model class, metric.

\section{Theory}
\label{sec:theory}

% TODO: ICL-as-GD, linear attention exactness, conditions for full transformers.

\section{Experiments}
\label{sec:experiments}

% TODO

\section{Related Work}
\label{sec:related}

% TODO: ICL-as-GD line; ROME/MEMIT; hypernetworks; distillation.

\section{Discussion}
\label{sec:discussion}

% TODO

\bibliographystyle{plainnat}
\bibliography{references}

\newpage
\input{checklist}

\end{document}
